% Options for packages loaded elsewhere
% Options for packages loaded elsewhere
\PassOptionsToPackage{unicode}{hyperref}
\PassOptionsToPackage{hyphens}{url}
%
\documentclass[
  english,
  russian,
  12pt,
  a4paper,
  DIV=11,
  numbers=noendperiod]{scrreprt}
\usepackage{xcolor}
\usepackage{amsmath,amssymb}
\setcounter{secnumdepth}{5}
\usepackage{iftex}
\ifPDFTeX
  \usepackage[T1]{fontenc}
  \usepackage[utf8]{inputenc}
  \usepackage{textcomp} % provide euro and other symbols
\else % if luatex or xetex
  \usepackage{unicode-math} % this also loads fontspec
  \defaultfontfeatures{Scale=MatchLowercase}
  \defaultfontfeatures[\rmfamily]{Ligatures=TeX,Scale=1}
\fi
\usepackage{lmodern}
\ifPDFTeX\else
  % xetex/luatex font selection
\fi
% Use upquote if available, for straight quotes in verbatim environments
\IfFileExists{upquote.sty}{\usepackage{upquote}}{}
\IfFileExists{microtype.sty}{% use microtype if available
  \usepackage[]{microtype}
  \UseMicrotypeSet[protrusion]{basicmath} % disable protrusion for tt fonts
}{}
\usepackage{setspace}
% Make \paragraph and \subparagraph free-standing
\makeatletter
\ifx\paragraph\undefined\else
  \let\oldparagraph\paragraph
  \renewcommand{\paragraph}{
    \@ifstar
      \xxxParagraphStar
      \xxxParagraphNoStar
  }
  \newcommand{\xxxParagraphStar}[1]{\oldparagraph*{#1}\mbox{}}
  \newcommand{\xxxParagraphNoStar}[1]{\oldparagraph{#1}\mbox{}}
\fi
\ifx\subparagraph\undefined\else
  \let\oldsubparagraph\subparagraph
  \renewcommand{\subparagraph}{
    \@ifstar
      \xxxSubParagraphStar
      \xxxSubParagraphNoStar
  }
  \newcommand{\xxxSubParagraphStar}[1]{\oldsubparagraph*{#1}\mbox{}}
  \newcommand{\xxxSubParagraphNoStar}[1]{\oldsubparagraph{#1}\mbox{}}
\fi
\makeatother


\usepackage{longtable,booktabs,array}
\usepackage{calc} % for calculating minipage widths
% Correct order of tables after \paragraph or \subparagraph
\usepackage{etoolbox}
\makeatletter
\patchcmd\longtable{\par}{\if@noskipsec\mbox{}\fi\par}{}{}
\makeatother
% Allow footnotes in longtable head/foot
\IfFileExists{footnotehyper.sty}{\usepackage{footnotehyper}}{\usepackage{footnote}}
\makesavenoteenv{longtable}
\usepackage{graphicx}
\makeatletter
\newsavebox\pandoc@box
\newcommand*\pandocbounded[1]{% scales image to fit in text height/width
  \sbox\pandoc@box{#1}%
  \Gscale@div\@tempa{\textheight}{\dimexpr\ht\pandoc@box+\dp\pandoc@box\relax}%
  \Gscale@div\@tempb{\linewidth}{\wd\pandoc@box}%
  \ifdim\@tempb\p@<\@tempa\p@\let\@tempa\@tempb\fi% select the smaller of both
  \ifdim\@tempa\p@<\p@\scalebox{\@tempa}{\usebox\pandoc@box}%
  \else\usebox{\pandoc@box}%
  \fi%
}
% Set default figure placement to htbp
\def\fps@figure{htbp}
\makeatother



\ifLuaTeX
\usepackage[bidi=basic,provide=*]{babel}
\else
\usepackage[bidi=default,provide=*]{babel}
\fi
% get rid of language-specific shorthands (see #6817):
\let\LanguageShortHands\languageshorthands
\def\languageshorthands#1{}


\setlength{\emergencystretch}{3em} % prevent overfull lines

\providecommand{\tightlist}{%
  \setlength{\itemsep}{0pt}\setlength{\parskip}{0pt}}



 
\usepackage[backend=biber,langhook=extras,autolang=other*]{biblatex}
\addbibresource{bib/cite.bib}

\usepackage[]{csquotes}

\usepackage{indentfirst}
\usepackage{float}
\floatplacement{figure}{H}
\IfFileExists{plex-otf.sty}{
  %% Full TeXlive
  \IfPackageAtLeastTF{plex-otf.sty}{2024-12-06}{
    \usepackage[math,RM={Scale=0.94},SS={Scale=0.94},SScon={Scale=0.94},TT={Scale=MatchLowercase,FakeStretch=0.9},DefaultFeatures={Ligatures=Common}]{plex-otf}
  }{
    \usepackage[RM={Scale=0.94},SS={Scale=0.94},SScon={Scale=0.94},TT={Scale=MatchLowercase,FakeStretch=0.9},DefaultFeatures={Ligatures=Common}]{plex-otf}
  }
}{
  %% TinyTeX
  \usepackage{libertine}
}
\KOMAoption{captions}{tableheading}
\\usepackage{indentfirst}
\makeatletter
\@ifpackageloaded{caption}{}{\usepackage{caption}}
\AtBeginDocument{%
\ifdefined\contentsname
  \renewcommand*\contentsname{Содержание}
\else
  \newcommand\contentsname{Содержание}
\fi
\ifdefined\listfigurename
  \renewcommand*\listfigurename{Список иллюстраций}
\else
  \newcommand\listfigurename{Список иллюстраций}
\fi
\ifdefined\listtablename
  \renewcommand*\listtablename{Список таблиц}
\else
  \newcommand\listtablename{Список таблиц}
\fi
\ifdefined\figurename
  \renewcommand*\figurename{Рисунок}
\else
  \newcommand\figurename{Рисунок}
\fi
\ifdefined\tablename
  \renewcommand*\tablename{Таблица}
\else
  \newcommand\tablename{Таблица}
\fi
}
\@ifpackageloaded{float}{}{\usepackage{float}}
\floatstyle{ruled}
\@ifundefined{c@chapter}{\newfloat{codelisting}{h}{lop}}{\newfloat{codelisting}{h}{lop}[chapter]}
\floatname{codelisting}{Список}
\newcommand*\listoflistings{\listof{codelisting}{Листинги}}
\makeatother
\makeatletter
\makeatother
\makeatletter
\@ifpackageloaded{caption}{}{\usepackage{caption}}
\@ifpackageloaded{subcaption}{}{\usepackage{subcaption}}
\makeatother
\usepackage{bookmark}
\IfFileExists{xurl.sty}{\usepackage{xurl}}{} % add URL line breaks if available
\urlstyle{same}
\hypersetup{
  pdftitle={Отчёт о лабораторной работе},
  pdfauthor={Бабаев Рахим},
  pdflang={ru-RU},
  hidelinks,
  pdfcreator={LaTeX via pandoc}}


\title{Отчёт о лабораторной работе}
\usepackage{etoolbox}
\makeatletter
\providecommand{\subtitle}[1]{% add subtitle to \maketitle
  \apptocmd{\@title}{\par {\large #1 \par}}{}{}
}
\makeatother
\subtitle{Лабораторная работа №5. Простые сети в GNS3. Анализ трафика}
\author{Бабаев Рахим}
\date{}
\begin{document}
\maketitle

\renewcommand*\contentsname{Содержание}
{
\setcounter{tocdepth}{1}
\tableofcontents
}
\listoffigures
\listoftables

\setstretch{1.5}
\chapter{Цель
работы}\label{ux446ux435ux43bux44c-ux440ux430ux431ux43eux442ux44b}

Целью лабораторной работы является построение простейших сетей в GNS3 на
базе коммутатора, а также маршрутизаторов FRR и VyOS, и анализ сетевого
трафика с помощью Wireshark.

\chapter{Задание}\label{ux437ux430ux434ux430ux43dux438ux435}

В ходе лабораторной работы требовалось выполнить следующие задания:

\begin{enumerate}
\def\labelenumi{\arabic{enumi}.}
\tightlist
\item
  Построить в GNS3 простейшую сеть на базе коммутатора Ethernet и двух
  оконечных устройств VPCS.
\item
  Назначить оконечным устройствам IP-адреса из сети 192.168.1.0/24 и
  проверить связь между узлами.
\item
  Запустить захват трафика и проанализировать ARP-, ICMP-, UDP- и
  TCP-сообщения в Wireshark.
\item
  Построить простейшую сеть на базе маршрутизатора FRR, коммутатора и
  оконечного устройства, выполнить настройку IP-адресации и проверить
  доступность маршрутизатора.
\item
  Построить простейшую сеть на базе маршрутизатора VyOS, выполнить его
  первоначальную настройку, назначить IP-адрес интерфейсу и проверить
  связь с оконечным устройством.
\end{enumerate}

\chapter{Теоретическое
введение}\label{ux442ux435ux43eux440ux435ux442ux438ux447ux435ux441ux43aux43eux435-ux432ux432ux435ux434ux435ux43dux438ux435}

GNS3 представляет собой платформу для моделирования и эмуляции сетей,
которая позволяет собирать виртуальные топологии из маршрутизаторов,
коммутаторов и оконечных узлов. Такой подход позволяет проверять работу
протоколов и взаимодействие устройств без использования физического
сетевого оборудования.

В данной лабораторной работе используются VPCS, Ethernet-коммутатор,
маршрутизаторы FRR и VyOS. VPCS применяется как лёгкое виртуальное
конечное устройство, на котором можно задавать IP-адрес, шлюз и
выполнять базовые сетевые команды. Коммутатор Ethernet обеспечивает
обмен кадрами в пределах локальной сети, а маршрутизаторы FRR и VyOS
используются для практической настройки интерфейсов и проверки сетевой
связности.

Для анализа сетевого взаимодействия используется Wireshark. С его
помощью можно захватывать трафик на выбранном соединении и анализировать
содержимое кадров и пакетов разных протоколов. В лабораторной работе
особое внимание уделяется ARP, ICMP, а также вариантам эхо-запросов в
режимах UDP и TCP.

Протокол ARP используется для определения MAC-адреса по известному
IPv4-адресу в пределах локального сегмента сети. Протокол ICMP
применяется для диагностических сообщений, в том числе для команды ping.
Анализ этих протоколов позволяет проследить, как на практике происходит
разрешение адресов, формирование эхо-запросов и ответов, а также
передача кадров между узлами через коммутатор и маршрутизатор.

FRR является программным маршрутизатором с интерфейсом командной строки,
близким по логике к оборудованию Cisco. VyOS также является программным
маршрутизатором, но использует собственную систему конфигурирования с
разделением на рабочий режим и режим настройки. Практическая работа с
обоими типами маршрутизаторов позволяет сравнить подходы к
конфигурированию сетевых устройств.

\chapter{Выполнение лабораторной
работы}\label{ux432ux44bux43fux43eux43bux43dux435ux43dux438ux435-ux43bux430ux431ux43eux440ux430ux442ux43eux440ux43dux43eux439-ux440ux430ux431ux43eux442ux44b}

\section{1. Моделирование простейшей сети на базе
коммутатора}\label{ux43cux43eux434ux435ux43bux438ux440ux43eux432ux430ux43dux438ux435-ux43fux440ux43eux441ux442ux435ux439ux448ux435ux439-ux441ux435ux442ux438-ux43dux430-ux431ux430ux437ux435-ux43aux43eux43cux43cux443ux442ux430ux442ux43eux440ux430}

\begin{enumerate}
\def\labelenumi{\arabic{enumi}.}
\item
  Сначала были запущены GNS3 VM и клиент GNS3, после чего был создан
  новый проект. Это позволило подготовить рабочую область для сборки
  топологии и дальнейшей проверки сетевого взаимодействия.
\item
  В рабочее пространство были добавлены коммутатор Ethernet и два узла
  VPCS. После этого устройствам были заданы осмысленные имена по шаблону
  из методических указаний, а сами узлы были соединены с коммутатором,
  чтобы получить простейшую локальную сеть. fileciteturn14file0
\item
  Затем была выполнена настройка IP-адресации для первого узла VPCS. Для
  PC-1 был задан адрес 192.168.1.11/24 и шлюз 192.168.1.1, после чего
  конфигурация была сохранена командой \texttt{save}.
  fileciteturn14file0
\item
  Аналогичным образом был настроен второй узел VPCS. Для PC-2 был задан
  адрес 192.168.1.12/24 в той же сети, что обеспечило возможность
  прямого обмена пакетами через коммутатор без необходимости
  маршрутизации. fileciteturn14file0
\item
  После назначения адресов была выполнена проверка связности. С помощью
  команды \texttt{ping} с одного узла был отправлен эхо-запрос на
  IP-адрес второго узла, и успешное получение ответа показало, что
  коммутация в локальном сегменте работает корректно.
  fileciteturn14file0
\item
  После завершения проверки работы сети все узлы проекта были
  остановлены. Это позволило завершить первый этап лабораторной работы и
  перейти к анализу трафика.
\end{enumerate}

\section{2. Анализ трафика в GNS3 посредством
Wireshark}\label{ux430ux43dux430ux43bux438ux437-ux442ux440ux430ux444ux438ux43aux430-ux432-gns3-ux43fux43eux441ux440ux435ux434ux441ux442ux432ux43eux43c-wireshark}

\begin{enumerate}
\def\labelenumi{\arabic{enumi}.}
\item
  Для анализа сетевого взаимодействия на соединении между PC-1 и
  коммутатором был запущен захват трафика. После выбора команды
  \texttt{Start\ capture} автоматически открылся Wireshark, а в
  топологии GNS3 на соответствующем соединении появился индикатор
  активного захвата. fileciteturn14file0
\item
  Затем в проекте были запущены все узлы. Сразу после старта в окне
  Wireshark появились кадры ARP, что связано с попытками устройств
  определить MAC-адреса по известным IP-адресам внутри локальной сети.
  fileciteturn14file0
\item
  При анализе ARP-трафика было установлено, что сначала формируется
  широковещательный ARP Request с вопросом о том, какому MAC-адресу
  соответствует определённый IPv4-адрес. В ответ узел с искомым адресом
  отправляет ARP Reply, после чего инициатор получает нужное
  соответствие IP-MAC и может передавать обычные IP-пакеты.
\item
  Далее в терминале PC-2 была просмотрена справка по команде
  \texttt{ping}, после чего был отправлен один эхо-запрос в ICMP-режиме
  на PC-1. В Wireshark были зафиксированы ICMP Echo Request и Echo
  Reply, что подтвердило корректную передачу диагностических сообщений в
  пределах локальной сети. fileciteturn14file0
\item
  После этого был выполнен один эхо-запрос в UDP-режиме. Анализ
  захваченного трафика показал, что в этом режиме пакет уже не является
  классическим ICMP-эхо, а передаётся как UDP-датаграмма, что позволяет
  проследить различие между типами диагностических сообщений и
  транспортных протоколов. fileciteturn14file0
\item
  Затем был выполнен один эхо-запрос в TCP-режиме. В Wireshark можно
  было увидеть TCP-сегменты, а также особенности установления
  соединения, характерные для данного протокола, включая обмен
  служебными флагами. fileciteturn14file0
\item
  На завершающем этапе захват трафика был остановлен. Полученные данные
  показали, что Wireshark позволяет на практике увидеть
  последовательность ARP-разрешения адресов и различия между ICMP-, UDP-
  и TCP-обменом.
\end{enumerate}

\section{3. Моделирование простейшей сети на базе маршрутизатора
FRR}\label{ux43cux43eux434ux435ux43bux438ux440ux43eux432ux430ux43dux438ux435-ux43fux440ux43eux441ux442ux435ux439ux448ux435ux439-ux441ux435ux442ux438-ux43dux430-ux431ux430ux437ux435-ux43cux430ux440ux448ux440ux443ux442ux438ux437ux430ux442ux43eux440ux430-frr}

\begin{enumerate}
\def\labelenumi{\arabic{enumi}.}
\item
  Для следующего этапа был создан новый проект в GNS3. В рабочее
  пространство были добавлены узел VPCS, коммутатор Ethernet и
  маршрутизатор FRR, после чего элементы были соединены в единую
  топологию. fileciteturn14file0
\item
  После размещения устройств им были назначены понятные имена в
  соответствии с методическими указаниями. Такой подход позволяет легче
  ориентироваться в топологии и отличать коммутатор, маршрутизатор и
  оконечный узел при дальнейшем анализе.
\item
  Далее на соединении между коммутатором и маршрутизатором был включён
  захват трафика. Это было необходимо для последующего анализа кадров и
  пакетов, проходящих между оконечным устройством и маршрутизатором FRR.
  fileciteturn14file0
\item
  После запуска всех устройств была открыта консоль каждого узла. На PC1
  был настроен IP-адрес 192.168.1.10/24 с указанием шлюза 192.168.1.1,
  затем конфигурация была сохранена и проверена командой
  \texttt{show\ ip}. fileciteturn14file0
\item
  Затем была выполнена настройка самого маршрутизатора FRR. Сначала ему
  было присвоено имя узла, затем на интерфейсе \texttt{eth0} был
  назначен адрес 192.168.1.1/24 и интерфейс был переведён в рабочее
  состояние командой \texttt{no\ shutdown}. После этого конфигурация
  была записана в память командой \texttt{write\ memory}.
  fileciteturn14file0
\item
  После завершения базовой настройки была выполнена проверка
  конфигурации маршрутизатора. Команды \texttt{show\ running-config} и
  \texttt{show\ interface\ brief} позволили убедиться, что интерфейс
  имеет правильный IP-адрес и находится в активном состоянии.
  fileciteturn14file0
\item
  Далее была выполнена проверка связности между оконечным устройством и
  маршрутизатором. Узел PC1 успешно отправлял эхо-запросы на адрес
  192.168.1.1, что означало корректную настройку локальной сети и
  работоспособность FRR как маршрутизатора в данной топологии.
  fileciteturn14file0
\item
  В Wireshark был проанализирован трафик, проходящий между коммутатором
  и маршрутизатором. В захваченном трафике были видны ARP-кадры для
  разрешения MAC-адресов и ICMP-сообщения, возникающие при проверке
  доступности маршрутизатора с узла PC1. fileciteturn14file0
\item
  После завершения анализа захват был остановлен, а все узлы в проекте
  были выключены. Это завершило этап работы с маршрутизатором FRR.
\end{enumerate}

\section{4. Моделирование простейшей сети на базе маршрутизатора
VyOS}\label{ux43cux43eux434ux435ux43bux438ux440ux43eux432ux430ux43dux438ux435-ux43fux440ux43eux441ux442ux435ux439ux448ux435ux439-ux441ux435ux442ux438-ux43dux430-ux431ux430ux437ux435-ux43cux430ux440ux448ux440ux443ux442ux438ux437ux430ux442ux43eux440ux430-vyos}

\begin{enumerate}
\def\labelenumi{\arabic{enumi}.}
\item
  Для последнего этапа лабораторной работы был создан отдельный проект,
  в который были добавлены VPCS, Ethernet-коммутатор и маршрутизатор
  VyOS. После размещения устройства были соединены, а их названия
  изменены по тому же принципу, что и в предыдущем задании.
  fileciteturn14file0
\item
  Затем на соединении между коммутатором и маршрутизатором был включён
  захват трафика. Это позволило подготовить топологию к анализу сетевого
  взаимодействия после настройки VyOS. fileciteturn14file0
\item
  После запуска всех устройств была выполнена настройка оконечного узла
  PC1. Ему был назначен IP-адрес 192.168.1.10/24 и шлюз 192.168.1.1,
  после чего конфигурация была сохранена и проверена.
  fileciteturn14file0
\item
  Далее была выполнена первоначальная настройка маршрутизатора VyOS.
  После загрузки был выполнен вход под учётной записью \texttt{vyos},
  при необходимости система была установлена на диск с помощью команды
  \texttt{install\ image}, а после завершения установки маршрутизатор
  был перезагружен. fileciteturn14file0
\item
  После перезагрузки был выполнен переход в режим конфигурирования
  командой \texttt{configure}. Затем маршрутизатору было присвоено имя
  узла, а интерфейсу \texttt{eth0} был назначен IP-адрес 192.168.1.1/24.
  fileciteturn14file0
\item
  Перед применением параметров была просмотрена разница конфигурации
  командой \texttt{compare}. После этого настройки были зафиксированы
  командой \texttt{commit} и сохранены командой \texttt{save}, что
  обеспечило их сохранение после возможной перезагрузки устройства.
  fileciteturn14file0
\item
  Далее была просмотрена информация об интерфейсах маршрутизатора
  командой \texttt{show\ interfaces}. Это позволило убедиться, что
  интерфейс \texttt{eth0} настроен правильно и доступен для обмена
  трафиком с узлом PC1. fileciteturn14file0
\item
  После завершения настройки была выполнена проверка подключений. Узел
  PC1 успешно отправлял эхо-запросы на адрес 192.168.1.1, что
  подтвердило корректную настройку сети и работоспособность
  маршрутизатора VyOS в собранной топологии. fileciteturn14file0
\item
  В окне Wireshark был проанализирован захваченный трафик. Как и в
  случае с FRR, в нём присутствовали ARP-кадры, обеспечивающие
  разрешение адресов, и ICMP-сообщения, отражающие процесс проверки
  доступности маршрутизатора. fileciteturn14file0
\item
  После завершения анализа захват был остановлен, все устройства были
  выключены, а работа с GNS3 завершена. Это означало полное выполнение
  всех заданий лабораторной работы.
\end{enumerate}

\chapter{Выводы}\label{ux432ux44bux432ux43eux434ux44b}

В ходе лабораторной работы были построены и исследованы несколько
простых сетевых топологий в GNS3. Были рассмотрены варианты сети на базе
коммутатора Ethernet, а также сети с использованием маршрутизаторов FRR
и VyOS.

Практическая часть показала, что даже простая локальная сеть включает
несколько уровней взаимодействия. До начала обмена IP-пакетами узлы
должны выполнить ARP-разрешение адресов, а уже затем становится
возможной передача ICMP-, UDP- и TCP-трафика. Анализ Wireshark позволил
наглядно проследить эту последовательность.

Отдельно была отработана базовая настройка программных маршрутизаторов
FRR и VyOS. Несмотря на различие в системах команд, оба маршрутизатора
позволяют назначать IP-адреса интерфейсам, сохранять конфигурацию и
обеспечивать сетевую связность между оконечным устройством и шлюзом.

Таким образом, лабораторная работа позволила закрепить практические
навыки построения простейших сетей в GNS3, захвата и анализа трафика в
Wireshark, а также базовой настройки программных маршрутизаторов в
виртуальной среде.

\chapter{Список
литературы}\label{ux441ux43fux438ux441ux43eux43a-ux43bux438ux442ux435ux440ux430ux442ux443ux440ux44b}

\begin{enumerate}
\def\labelenumi{\arabic{enumi}.}
\tightlist
\item
  Королькова А. В., Кулябов Д. С. Сетевые технологии. Лабораторный
  практикум.
\item
  Документация Wireshark.
\item
  Документация GNS3.
\item
  Документация FRRouting.
\item
  Документация VyOS.
\end{enumerate}


\printbibliography



\end{document}
